\def\sectiontag{Patterns are mathematical information, not decoration}
\section{Special matrices and visible structure}

\begin{frame}{Structure changes the calculation}
  \begin{columns}[T,onlytextwidth]
    \begin{column}{0.46\textwidth}
      \[
      \begin{pmatrix}
      4&0&0\\
      0&-2&0\\
      0&0&5
      \end{pmatrix}
      \]
    \end{column}
    \begin{column}{0.50\textwidth}
      \begin{block}{Visible immediately}
        no off-diagonal interaction
      \end{block}
      \begin{exampleblock}{Consequence}
        powers, products and interpretation become transparent
      \end{exampleblock}
    \end{column}
  \end{columns}

  \vspace{0.5cm}
  \concept{Seeing the pattern is part of solving the problem.}
\end{frame}

\begin{frame}{Zero and identity: two neutral objects}
  \begin{columns}[T,onlytextwidth]
    \begin{column}{0.46\textwidth}
      \begin{block}{Zero matrix}
        \[
        A+0=A
        \]
        neutral for addition
      \end{block}
    \end{column}
    \begin{column}{0.46\textwidth}
      \begin{exampleblock}{Identity matrix}
        \[
        I_mA=A,\qquad AI_n=A
        \]
        neutral for multiplication
      \end{exampleblock}
    \end{column}
  \end{columns}
\end{frame}

\begin{frame}{Identity sizes must fit}
  \begin{columns}[T,onlytextwidth]
    \begin{column}{0.48\textwidth}
      \[
      A\in\R^{m\times n}
      \]
      \begin{block}{Left identity}
        \[
        I_mA=A
        \]
        $I_m$ has shape $m\times m$
      \end{block}
    \end{column}
    \begin{column}{0.48\textwidth}
      \vspace{0.35cm}
      \begin{exampleblock}{Right identity}
        \[
        AI_n=A
        \]
        $I_n$ has shape $n\times n$
      \end{exampleblock}
    \end{column}
  \end{columns}

  \vspace{0.45cm}
  \takeaway{The identity is conceptually one object, but dimensionally many matrices.}
\end{frame}

\begin{frame}{Transpose: rows become columns}
  \begin{columns}[T,onlytextwidth]
    \begin{column}{0.40\textwidth}
      \[
      A=\begin{pmatrix}
      1&2&3\\
      4&5&6
      \end{pmatrix}
      \]
      \centering $2\times3$
    \end{column}
    \begin{column}{0.06\textwidth}
      \centering\vspace{0.85cm}\Large $\longrightarrow$
    \end{column}
    \begin{column}{0.40\textwidth}
      \[
      A^{\mathsf T}=\begin{pmatrix}
      1&4\\
      2&5\\
      3&6
      \end{pmatrix}
      \]
      \centering $3\times2$
    \end{column}
  \end{columns}
\end{frame}

\begin{frame}{Transpose rules and symmetry}
  \begin{columns}[T,onlytextwidth]
    \begin{column}{0.48\textwidth}
      \[
      (A^{\mathsf T})^{\mathsf T}=A
      \]
      \[
      (A+B)^{\mathsf T}=A^{\mathsf T}+B^{\mathsf T}
      \]
      \[
      (AB)^{\mathsf T}=B^{\mathsf T}A^{\mathsf T}
      \]
    \end{column}
    \begin{column}{0.48\textwidth}
      \begin{exampleblock}{Symmetric matrix}
        \[
        A^{\mathsf T}=A
        \]
        entries mirror across the main diagonal
      \end{exampleblock}
    \end{column}
  \end{columns}
\end{frame}
